\section{Анализ предметной области}

\subsection{Особенности работы йога-студий и их деятельности}

Йога как практика физического и духовного развития существует тысячи лет. В современном мире йога стала массовым явлением: миллионы людей занимаются ею для поддержания здоровья, снятия стресса и улучшения качества жизни. Йога-студии предлагают различные направления, включая хатха-йогу, винтьясу, аштангу, кундалини, инь-йогу и медитацию.

Современные йога-студии сталкиваются с рядом проблем. В первую очередь, это необходимость управления расписанием занятий и обработка записей клиентов. Также требуется информирование клиентов об изменениях в расписании, отслеживание прогресса клиентов и сбор обратной связи в виде отзывов и рейтингов. Кроме того, важной задачей является управление информацией об инструкторах.

\subsection{Проблематика традиционного подхода к управлению записями}

Традиционные методы управления, такие как запись по телефону, ведение журналов или таблиц Excel, являются неэффективными и отнимают много времени как у администраторов, так и у клиентов. При отсутствии специализированной системы данные о клиентах и записях распределены по разным источникам, что приводит к ряду типовых проблем.

Разрозненность данных проявляется в том, что одни и те же записи дублируются в нескольких местах, создавая расхождения в информации о клиентах и занятиях. Сложность актуализации заключается в том, что любое изменение расписания требует ручной правки в нескольких каналах, из-за чего часть информации быстро устаревает. Высокие издержки пользовательского поиска вынуждают клиента переключаться между источниками, чтобы получить полную картину о занятиях. Отсутствие единого профиля пользователя не позволяет отслеживать историю посещений и прогресс. Кроме того, клиенты не имеют возможности получать своевременные напоминания о предстоящих занятиях.

\subsection{Обоснование необходимости автоматизации процессов}

Автоматизация в рассматриваемой предметной области необходима не сама по себе, а как инструмент устранения перечисленных противоречий. Для йога-студии автоматизация обеспечивает единое информационное пространство, в котором данные о клиентах, расписании, записях и отзывах поддерживаются согласованно.

Практический эффект автоматизации проявляется в нескольких аспектах. Сокращается время обновления расписания и обработки записей клиентов. Снижается вероятность ошибок при повторном ручном вводе данных. Повышается доступность информации для конечного пользователя. Упрощается отслеживание прогресса клиентов. Формируется единая логика управления профилями и записями. Таким образом, автоматизация является необходимым условием устойчивой работы сервиса, ориентированного на регулярное использование и постоянное развитие.

\subsection{Структура предметной области и ключевые сущности}

На основании анализа деятельности сервиса выделяются следующие ключевые сущности. Пользователь представляет собой клиента йога-студии, имеющего профиль с личными данными, уровнем подготовки и историей записей. Занятие является основной единицей, включающей название, описание, время проведения, инструктора, город и максимальное количество участников. Инструктор — это специалист, проводящий занятия, с информацией об имени, специализации и городе работы. Запись представляет собой связь между пользователем и занятием, фиксирующую намерение клиента посетить занятие. Отзыв содержит оценку и комментарий пользователя о посещённом занятии. Прогресс отражает статистику посещений пользователя за определённый период.

\subsubsection{Детальное описание сущности «Пользователь»}

Пользователь является центральной сущностью системы. Каждый пользователь имеет уникальный идентификатор, адрес электронной почты (используемый для входа), пароль (хранится в зашифрованном виде), имя, отображаемое в профиле, город проживания (влияет на фильтрацию занятий), уровень подготовки (новичок, средний, продвинутый), а также аватар (опционально). Пользователь может просматривать каталог занятий, записываться на них, отменять запись, оставлять отзывы и отслеживать свой прогресс.

\subsubsection{Детальное описание сущности «Занятие»}

Занятие представляет собой основную единицу, предоставляемую йога-студией. Каждое занятие характеризуется следующими атрибутами: уникальный идентификатор, название (например, «Хатха-йога для начинающих»), подробное описание, день и время проведения, продолжительность, инструктор, город проведения, максимальное количество участников, текущее количество записанных, цена (если занятие платное), а также флаг публикации (видимо для пользователей или скрыто). Занятия группируются по направлениям и городам.

\subsubsection{Детальное описание сущности «Инструктор»}

Инструктор — это специалист, проводящий занятия. Каждый инструктор имеет уникальный идентификатор, имя, специализацию (направления йоги, которые он преподаёт), город работы, краткую биографию, фотографию и рейтинг (рассчитываемый на основе отзывов пользователей). Инструктор может вести несколько занятий в разных городах. Пользователи могут просматривать профиль инструктора и оставлять отзывы о его занятиях.

\subsubsection{Детальное описание сущности «Запись»}

Запись фиксирует намерение пользователя посетить конкретное занятие. Запись содержит идентификатор пользователя, идентификатор занятия, дату и время создания записи, статус (активна, отменена, посещена). Система проверяет наличие свободных мест перед созданием записи. При отмене записи место освобождается и становится доступным для других пользователей. После проведения занятия статус записи меняется на «посещено», и пользователь получает возможность оставить отзыв.

\subsubsection{Детальное описание сущности «Отзыв»}

Отзыв позволяет пользователю оценить качество проведённого занятия. Отзыв содержит идентификатор пользователя, идентификатор занятия, рейтинг (от 1 до 5 звёзд), текстовый комментарий, дату создания. Пользователь может оставить отзыв только после фактического посещения занятия (статус записи «посещено»). Отзывы влияют на рейтинг инструктора и помогают другим пользователям при выборе занятий.

\subsubsection{Детальное описание сущности «Прогресс»}

Прогресс отражает статистику посещений пользователя за определённый период. Система автоматически обновляет прогресс после каждого посещённого занятия. Прогресс включает общее количество занятий, количество занятий по месяцам, количество занятий по направлениям йоги, а также динамику изменения уровня подготовки. Данные о прогрессе визуализируются в виде графиков с использованием библиотеки MPAndroidChart.

\subsection{Анализ существующих решений в тематике записи на занятия}

На рынке представлено несколько типов решений для записи на занятия йогой. В таблице \ref{tab:comparison} приведён сравнительный анализ существующих аналогов.

\vspace{1em}
\renewcommand{\arraystretch}{1.4}
\begin{table}[H]
	\centering
	\caption{Сравнительный анализ существующих решений}
	\label{tab:comparison}
	\normalsize
	\setlength{\tabcolsep}{6pt}
	\begin{tabular}{|m{4.5cm}|m{2.2cm}|m{2.2cm}|m{2.9cm}|m{2.8cm}|}
		\hline
		Критерий & TimePad & FitStars & Excel/ручной & Наше решение \\
		\hline
		Бесплатно & Нет & Нет & Да & Да \\
		\hline
		Мобильное приложение & Нет & Да & Нет & Да \\
		\hline
		Запись онлайн & Да & Да & Нет & Да \\
		\hline
		Профиль пользователя & Нет & Да & Нет & Да \\
		\hline
		Отслеживание прогресса & Нет & Частично & Нет & Да \\
		\hline
		Отзывы и рейтинг & Нет & Да & Нет & Да \\
		\hline
		Уведомления & Нет & Да & Нет & Да \\
		\hline
		Офлайн-режим & Нет & Нет & Да & Да \\
		\hline
	\end{tabular}
\end{table}
\renewcommand{\arraystretch}{1}

\subsubsection{Зарубежные платформы}

На международном рынке существует ряд решений для управления фитнес-студиями. Платформа Mindbody является одним из лидеров, предоставляя комплексные возможности для управления расписанием, записями и платежами. Однако данный сервис является платным и ориентирован преимущественно на крупные сети студий. Сервис Zen Planner предлагает аналогичный функционал, но также требует ежемесячной платы и не имеет полноценной поддержки русского языка. Ключевым недостатком зарубежных решений является их ориентация на западный рынок и отсутствие бесплатных тарифов для малых студий.

\textbf{Платформа Mindbody} — это комплексное решение для управления фитнес-бизнесом, включающее модули для расписания, записи клиентов, обработки платежей, маркетинга и аналитики. Система интегрируется с различными платёжными шлюзами и предлагает мобильное приложение для клиентов. Однако стоимость подписки начинается от 129 долларов в месяц, что недоступно для малых йога-студий.

\textbf{Платформа ClassPass} — это не столько инструмент управления студией, сколько агрегатор занятий, позволяющий пользователям покупать абонементы и записываться на занятия в различных студиях. С точки зрения студии, ClassPass привлекает новых клиентов, но студия теряет контроль над ценообразованием и взаимоотношениями с клиентами.

\subsubsection{Региональные и отечественные решения}

В Российской Федерации также представлены решения для управления занятиями. FitStars является мобильным приложением, предлагающим каталог занятий и запись онлайн, однако его функционал ограничен, а некоторые возможности требуют платной подписки. TimePad представляет собой платформу для управления мероприятиями, которая может использоваться для записи на занятия, но не имеет мобильного приложения и ориентирована скорее на разовые события, чем на регулярные занятия. Многие студии продолжают использовать ручную запись через Excel и мессенджеры, что подтверждает актуальность разработки нового решения.

\textbf{Платформа FitStars} предлагает бесплатную версию с ограниченным функционалом (до 50 записей в месяц). Платная подписка (от 990 рублей в месяц) открывает доступ к неограниченному количеству записей, статистике и интеграции с социальными сетями. Однако приложение не предоставляет возможности отслеживания прогресса клиентов и не имеет офлайн-режима.

\textbf{Платформа TimePad} изначально создавалась для управления мероприятиями (конференции, семинары, тренинги). Она поддерживает создание событий, продажу билетов и рассылку уведомлений. Однако интерфейс не оптимизирован для регулярных занятий (нет понятия «абонемент», «прогресс», «инструктор»), а мобильное приложение отсутствует.

\subsubsection{Выводы по анализу существующих решений}

Проведённый анализ позволяет выделить несколько ключевых принципов, на которых строятся успешные современные приложения для записи на занятия. Пользователи ожидают возможность записываться на занятия со смартфона в любое время, поэтому мобильная платформа является обязательным требованием. Приложение должно сохранять базовую функциональность даже без подключения к интернету, обеспечивая автономную работу. Наличие профиля пользователя, отслеживание прогресса и персональные рекомендации повышают вовлечённость. Возможность оставлять отзывы и оценки формирует доверие к сервису. Напоминания о занятиях снижают количество пропусков и повышают удовлетворённость клиентов.

Существующие решения, такие как TimePad и FitStars, либо не поддерживают мобильные устройства, либо являются платными, либо не предоставляют полного набора функций для клиентов йога-студий. Таким образом, разработка собственного мобильного приложения «Йога-студия» является актуальной и целесообразной.

\subsection{Анализ целевой аудитории}

Целевая аудитория разрабатываемого приложения может быть разделена на две основные группы: клиенты йога-студий и администраторы (включая инструкторов).

\subsubsection{Портрет клиента йога-студии}

Клиент йога-студии — это человек в возрасте от 18 до 55 лет, ведущий активный образ жизни, интересующийся здоровьем и саморазвитием. Основные потребности клиента: возможность быстро записаться на занятие, получать напоминания, отслеживать свой прогресс, оставлять отзывы. Клиент предпочитает использовать мобильное приложение, так как это удобно и быстро. Большинство клиентов используют смартфоны на базе Android (около 70\% рынка).

\subsubsection{Портрет администратора йога-студии}

Администратор йога-студии — это сотрудник студии, отвечающий за управление расписанием, обработку записей и взаимодействие с клиентами. Основные потребности администратора: возможность быстро изменять расписание, просматривать список записанных клиентов, отслеживать загруженность залов и инструкторов, получать отчёты.

\subsubsection{Портрет инструктора}

Инструктор — это специалист, проводящий занятия. Основные потребности инструктора: просмотр расписания своих занятий, просмотр списка записанных клиентов, возможность отмечать посещаемость, отслеживание своего рейтинга на основе отзывов.

\subsection{Выводы по разделу}

На основе проведённого анализа было принято решение разработать мобильное приложение «Йога-студия» на платформе Android с использованием языка Kotlin. Приложение будет включать полный цикл записи на занятия, профиль пользователя, статистику, отзывы и уведомления. Выбранная архитектура обеспечивает модульность, удобство поддержки и масштабирования.